% Einheit 1 von 5 – Tangentengleichung im Punkt (bank/tangente-normale-schnittwinkel/e1.jsonl, zusammenbau v0.1)
%% TODO Zeitmarke Einheit 1: Marken-Zeile nicht lesbar
%% TODO Zweigzeile Teil 1: was hier gelernt wird (Satz in Schülersprache aus dem Einheitstitel) – Einheit 1
\einheitenkopf[e1]{Einheit 1 von 5 · Tangentengleichung im Punkt}
\zweigzeile{Abitur GK}
\setcounter{aufgabe}{8}

%% TODO Titel: Ich-kann-Satz fehlt in der Bank (Platzhalter: Kettenname) – Nr. 9
%% TODO Anweisung: ein Satz über der ersten Teilaufgabe fehlt in der Bank – Nr. 9
\begin{aufgabe}{Tangentengleichung}
%% TODO \rechenplatz{n}: Schrittzahl des Grundfalls fehlt in der Bank (nur bei fester Schreibform, 2.3 b)
\begin{teile}
\teil Gesucht ist die Tangente an den Graphen von $f$ im Punkt $P(2 | f(2))$. Was ist gegeben? \\ \kreuz{ein Punkt} \\ \kreuz{eine Steigung} \\ \kreuz{ein Winkel}
\teil $f(x) = x^3 - 2x$ – Tangente im Punkt $(2 | f(2))$? $t(x) =$ \leerfeld
\teil Liegt $P(-2 | 9)$ auf dem Graphen von $f(x) = -x^3 + 2x^2 + 3x - 1$? Wenn ja: Anstieg und Tangente in $P$? \hfill (FHR 2020) \\ Anstieg: \leerfeld , $t(x) =$ \leerfeld
\teil Ein Gartenhaus hat ein symmetrisches Satteldach: Das Haus ist $6$ m breit, der First liegt $2$ m über der Traufe; $1$ LE entspricht $1$ m. Anstieg der linken und der rechten Dachschräge, und wie lang ist eine Dachschräge? \hfill (FHR 2021) \\ links: \leerfeld , rechts: \leerfeld , Länge: \leerfeld[m]
\teil $f(x) = 2e^x - 3$ – Gleichung der Tangente $t$ im Schnittpunkt des Graphen mit der $y$-Achse? \hfill (Abitur 2026 GK) Kontrolle: $y = 2x - 1$. $t(x) =$ \leerfeld
\teil $f(x) = x^3 - 3x^2 + 1$ – in welchen Punkten hat der Graph die Steigung $9$? \leerfeld
\teil $f(x) = (x^2 - 5) \cdot e^x$ mit $f'(x) = (x^2 + 2x - 5) \cdot e^x$: Weise nach, dass der Graph an der Stelle $x = -1 + \sqrt{6}$ eine waagerechte Tangente hat, und gib den Funktionswert dort an (auf zwei Stellen). \hfill (Abitur 2023 GK) \\ $f \approx$ \leerfeld
\teil Die Abbildung zeigt den Graphen von $f(x) = x^3 - 3x$. An welcher weiteren Stelle hat der Graph dieselbe Steigung wie bei $x = 2$? Zeichne die Tangente dort ein.

\begin{ksys}[xmin=-3,xmax=3,ymin=-5,ymax=5] \funktion{\x^3-3*\x}{f} \end{ksys}
$x =$ \leerfeld
\teil Die Abbildung zeigt den Graphen von $f$ und die Tangente $t$ im Punkt $(4 | f(4))$. Gleichung von $t$ anhand der Abbildung? \hfill (Abitur 2024 LK) \\

\begin{ksys}[xmin=-2,xmax=6,ymin=-5,ymax=6,ablesen] \funktion{0.25*\x^2}{f} \tangentean*{0.25*\x^2}{4}{t} \end{ksys}
$t(x) =$ \leerfeld
\teil Die Gerade $g$ verläuft durch die Punkte $(1 | f(1))$ und $(-2 | f(-2))$ des Graphen von $f(x) = x^4 - 3x^2$. Weise nach, dass $g$ die Tangente an den Graphen im Punkt $(1 | f(1))$ ist. \hfill (Abitur 2017 LK)
\teil Für $a > 0$ ist $f_a(x) = a \cdot x^3$. Weise nach, dass die Tangente an den Graphen von $f_a$ im Punkt $(u | f_a(u))$ die $y$-Achse im Punkt $(0 | -2f_a(u))$ schneidet. \hfill (Abitur 2024 LK)
\end{teile}
\end{aufgabe}

%% TODO Titel: Ich-kann-Satz fehlt in der Bank (Platzhalter: Kettenname) – Nr. 10
%% TODO Anweisung: ein Satz über der ersten Teilaufgabe fehlt in der Bank – Nr. 10
\begin{aufgabe}{Stellen, deren Tangente die y-Achse oberhalb einer Schranke schneidet, über eine Ungleichung bestimmen}
\begin{teile}
\teil $f(x) = \sqrt{x^2 + 16}$ – für welche Stellen $c$ schneidet die Tangente im Punkt $(c | f(c))$ die $y$-Achse oberhalb von $y = 2$? \hfill (Abitur 2018 LK) \\ \leerfeld
\end{teile}
\end{aufgabe}

%% TODO Titel: Ich-kann-Satz fehlt in der Bank (Platzhalter: Kettenname) – Nr. 11
%% TODO Anweisung: ein Satz über der ersten Teilaufgabe fehlt in der Bank – Nr. 11
\begin{aufgabe}{Parallelität der Wendetangenten einer Schar über eine parameterunabhängige Steigung nachweisen}
\begin{teile}
\teil Für $k > 0$ ist $f_k(x) = k^2x^3 - 3kx^2 + 4x$. Weise nach, dass die Wendetangenten aller Graphen zueinander parallel sind. \hfill (Abitur 2017 LK)
\end{teile}
\end{aufgabe}

%% TODO Titel: Ich-kann-Satz fehlt in der Bank (Platzhalter: Kettenname) – Nr. 12
%% TODO Anweisung: ein Satz über der ersten Teilaufgabe fehlt in der Bank – Nr. 12
\begin{aufgabe}{Tangentengleichung – Fehler finden, Begründen}
\begin{teile}
\teil Ich finde den Fehler: Ben bestimmt die Tangente an $f(x) = x^2 + 3x$ im Punkt $(1 | f(1))$: \rechnung{m &= f(1) = 4 \\ t(x) &= 4x}
\teil Begründe, warum die Tangente im Berührpunkt dieselbe Steigung hat wie der Graph.
\end{teile}
\end{aufgabe}
